\documentclass[10pt,a4paper]{amsart}
\usepackage{amsmath,amssymb,amsthm,graphicx}
% Retained minimal preamble for the Stacks Project chapter more-morphisms.tex (More on Morphisms).
% Verbatim excerpt of the chapter preamble: Xy-pic setup (native preamble lines 14-18) and the native
% theorem declarations with the mathematical macros (native preamble lines 154-217). Copied mechanically;
% nothing added, removed or redefined. The wrapper class and the amsmath/amssymb/amsthm/graphicx base
% packages are supplied by the assembler, so the class/externaldocument part of the native preamble is omitted.
% For commutative diagrams we use Xy-pic
\usepackage[all]{xy}

% We use 2cell for 2-commutative diagrams.
\xyoption{2cell}
% Theorem environments.
%
\theoremstyle{plain}
\newtheorem{theorem}[subsection]{Theorem}
\newtheorem{proposition}[subsection]{Proposition}
\newtheorem{lemma}[subsection]{Lemma}

\theoremstyle{definition}
\newtheorem{definition}[subsection]{Definition}
\newtheorem{example}[subsection]{Example}
\newtheorem{exercise}[subsection]{Exercise}
\newtheorem{situation}[subsection]{Situation}

\theoremstyle{remark}
\newtheorem{remark}[subsection]{Remark}
\newtheorem{remarks}[subsection]{Remarks}

\numberwithin{equation}{subsection}

% Macros
%
\def\lim{\mathop{\mathrm{lim}}\nolimits}
\def\colim{\mathop{\mathrm{colim}}\nolimits}
\def\Spec{\mathop{\mathrm{Spec}}}
\def\Hom{\mathop{\mathrm{Hom}}\nolimits}
\def\Ext{\mathop{\mathrm{Ext}}\nolimits}
\def\SheafHom{\mathop{\mathcal{H}\!\mathit{om}}\nolimits}
\def\SheafExt{\mathop{\mathcal{E}\!\mathit{xt}}\nolimits}
\def\Sch{\mathit{Sch}}
\def\Mor{\mathop{\mathrm{Mor}}\nolimits}
\def\Ob{\mathop{\mathrm{Ob}}\nolimits}
\def\Sh{\mathop{\mathit{Sh}}\nolimits}
\def\NL{\mathop{N\!L}\nolimits}
\def\CH{\mathop{\mathrm{CH}}\nolimits}
\def\proetale{{pro\text{-}\acute{e}tale}}
\def\etale{{\acute{e}tale}}
\def\QCoh{\mathit{QCoh}}
\def\Ker{\mathop{\mathrm{Ker}}}
\def\Im{\mathop{\mathrm{Im}}}
\def\Coker{\mathop{\mathrm{Coker}}}
\def\Coim{\mathop{\mathrm{Coim}}}

% Boxtimes
%
\DeclareMathSymbol{\boxtimes}{\mathbin}{AMSa}{"02}

%
% Macros for moduli stacks/spaces
%
\def\QCohstack{\mathcal{QC}\!\mathit{oh}}
\def\Cohstack{\mathcal{C}\!\mathit{oh}}
\def\Spacesstack{\mathcal{S}\!\mathit{paces}}
\def\Quotfunctor{\mathrm{Quot}}
\def\Hilbfunctor{\mathrm{Hilb}}
\def\Curvesstack{\mathcal{C}\!\mathit{urves}}
\def\Polarizedstack{\mathcal{P}\!\mathit{olarized}}
\def\Complexesstack{\mathcal{C}\!\mathit{omplexes}}
% \Pic is the operator that assigns to X its picard group, usage \Pic(X)
% \Picardstack_{X/B} denotes the Picard stack of X over B
% \Picardfunctor_{X/B} denotes the Picard functor of X over B
\def\Pic{\mathop{\mathrm{Pic}}\nolimits}
\def\Picardstack{\mathcal{P}\!\mathit{ic}}
\def\Picardfunctor{\mathrm{Pic}}
\def\Deformationcategory{\mathcal{D}\!\mathit{ef}}
\newenvironment{reference}{}{}
\newenvironment{slogan}{}{}
\newenvironment{history}{}{}
\usepackage{cleveref}
\crefname{originallemma}{original lemma}{original lemmas}
\Crefname{originallemma}{Original lemma}{Original lemmas}
\makeatletter
\newlabel{algebra-lemma-etale-makes-quasi-finite-finite-one-prime}{{00UJ}{source}{}{sourcepointer-1}{}}
\newlabel{algebra-lemma-etale-makes-quasi-finite-finite-one-prime@cref}{{[originallemma][0][]00UJ}{[1][0][]source}{}{}{}}
\makeatother
\crefname{originallemma}{original lemma}{original lemmas}
\Crefname{originallemma}{Original lemma}{Original lemmas}
\makeatletter
\newlabel{equation-basic-diagram}{{02LJ}{source}{}{sourcepointer-2}{}}
\newlabel{equation-basic-diagram@cref}{{[originallemma][0][]02LJ}{[1][0][]source}{}{}{}}
\makeatother
\crefname{originallemma}{original lemma}{original lemmas}
\Crefname{originallemma}{Original lemma}{Original lemmas}
\makeatletter
\newlabel{morphisms-lemma-quasi-finite-at-point-characterize}{{01TH}{source}{}{sourcepointer-3}{}}
\newlabel{morphisms-lemma-quasi-finite-at-point-characterize@cref}{{[originallemma][0][]01TH}{[1][0][]source}{}{}{}}
\makeatother
\title{An isolated point of a finite type fibre becomes finite over an elementary etale neighbourhood}
\author{The Stacks Project Authors}
\date{stacks-project:more-morphisms.tex:local-snapshot-20261011; GNU FDL 1.2 or later}
\begin{document}
\maketitle
\noindent\textit{Editorial note.} Statement span is lines 11923-11953 of sources/stacks-more-morphisms/pinned-original-source.tex; proof span is lines 11955-11963, the complete author proof (7 body lines).\par\medskip
\noindent\textit{Editorial note.} Navigation stubs: 3 cited-but-not-included labels are supplied as non-load-bearing pointers with their pinned tags where the snapshot tag index has them.\par\medskip
\noindent\textit{Editorial note.} extra\_packages is empty on purpose: the retained preamble already loads xy with options and a second declaration clashes in the assembler wrapper.\par\medskip
\noindent\textit{Editorial note.} Text is retained verbatim, including the author's own typos and placeholders.\par\medskip
\noindent\textit{Allowed standard prerequisites (editorial).}\begin{itemize}
\item Explicit prerequisite boundary. The characterization of quasi-finiteness at a point as an isolated fibre point under finite type (Morphisms, cited in the statement footnote) and the affine algebra case (Algebra, etale-makes-quasi-finite-finite-one-prime); cited inputs with the localization reduction written in the retained span.
\end{itemize}\medskip
\ifdefined\hypertarget\hypertarget{sourcepointer-1}{}\fi
\noindent\textit{Source reference.} Original lemma 00UJ: sources/stacks-more-morphisms/pinned-original-source.tex, source lines 11963--11963. Cited by the retained statement or proof but not part of the collected spans; supplied only as a navigation pointer. The argument does not depend on its content.\par\medskip
\ifdefined\hypertarget\hypertarget{sourcepointer-2}{}\fi
\noindent\textit{Source reference.} Original lemma 02LJ: sources/stacks-more-morphisms/pinned-original-source.tex, source lines 11909--11916. Cited by the retained statement or proof but not part of the collected spans; supplied only as a navigation pointer. The argument does not depend on its content.\par\medskip
\ifdefined\hypertarget\hypertarget{sourcepointer-3}{}\fi
\noindent\textit{Source reference.} Original lemma 01TH: sources/stacks-more-morphisms/pinned-original-source.tex, source lines 11963--11963. Cited by the retained statement or proof but not part of the collected spans; supplied only as a navigation pointer. The argument does not depend on its content.\par\medskip
\begin{lemma}
\label{lemma-etale-makes-quasi-finite-finite-at-point}
Let $f : X \to S$ be a morphism of schemes.
Let $x \in X$. Set $s = f(x)$.
Assume that
\begin{enumerate}
\item $f$ is locally of finite type, and
\item $x \in X_s$ is isolated\footnote{In the presence of (1)
this means that $f$ is
quasi-finite at $x$, see
Morphisms, Lemma \ref{morphisms-lemma-quasi-finite-at-point-characterize}.}.
\end{enumerate}
Then there exist
\begin{enumerate}
\item[(a)] an elementary \'etale neighbourhood $(U, u) \to (S, s)$,
\item[(b)] an open subscheme $V \subset X_U$
(see \ref{equation-basic-diagram})
\end{enumerate}
such that
\begin{enumerate}
\item[(\romannumeral1)] $V \to U$ is a finite morphism,
\item[(\romannumeral2)] there is a unique point $v$ of $V$
mapping to $u$ in $U$, and
\item[(\romannumeral3)] the point $v$ maps to $x$
under the morphism $X_U \to X$, inducing $\kappa(x) = \kappa(v)$.
\end{enumerate}
Moreover, for any elementary \'etale neighbourhood $(U', u') \to (U, u)$
setting $V' = U' \times_U V \subset X_{U'}$ the triple $(U', u', V')$
satisfies the properties
(\romannumeral1), (\romannumeral2), and (\romannumeral3) as well.
\end{lemma}

\begin{proof}
Let $Y \subset X$, $W \subset S$ be affine opens such that
$f(Y) \subset W$ and such that $x \in Y$. Note that $x$ is
also an isolated point of the fibre of the morphism $f|_Y : Y \to W$.
If we can prove the theorem for $f|_Y : Y \to W$, then the
theorem follows for $f$. Hence we reduce to the case where
$f$ is a morphism of affine schemes. This case is
Algebra, Lemma \ref{algebra-lemma-etale-makes-quasi-finite-finite-one-prime}.
\end{proof}

\end{document}
